\documentclass[../../../include/open-logic-section]{subfiles}

\begin{document}

\olfileid[id]{sth}{ord-arithmetic}{using-addition}
\olsection{Menggunakan Penjumlahan Ordinal}

Dengan menggunakan penjumlahan pada ordinal, kita dapat menghitung secara
eksplisit peringkat berbagai himpunan, dalam pengertian
\olref[spine][rank]{defnsetrank}:

\begin{lem}\ollabel{rankcomputation}
Jika $\setrank{A} = \alpha$ dan $\setrank{B} = \beta$, maka:
\begin{enumerate}
	\item\ollabel{exrankpow} $\setrank{\Pow{A}} = \alpha\ordplus 1$
	\item\ollabel{exrankpair} $\setrank{\{A, B\}} = \max(\alpha,
	\beta) \ordplus 1$
	\item\ollabel{exrankcup} $\setrank{A \cup B} = \max(\alpha,
	\beta)$
	\item\ollabel{exranktuple} $\setrank{\tuple{A,B}} = \max(\alpha,
	\beta) \ordplus  2$
	\item\ollabel{exranktimes} $\setrank{A \times B} \leq \max(\alpha,
	\beta) \ordplus  2$
	\item\ollabel{exrankunion} $\setrank{\bigcup A} = \alpha$ jika
	$\alpha$ kosong atau merupakan ordinal limit; $\setrank{\bigcup A} = \gamma$
	jika $\alpha = \gamma\ordplus 1$
\end{enumerate}
\end{lem}

\begin{proof}
Sepanjang bukti ini, kita berulang kali menggunakan
\olref[spine][rank]{ranksupstrict}.

\emph{\olref{exrankpow}.} Jika $x \subseteq A$, maka $\setrank{x} \leq
\setrank{A}$. Jadi $\setrank{\Pow{A}} \leq \alpha \ordplus  1$. Karena
khususnya $A \in \Pow{A}$, maka $\setrank{\Pow{A}} = \alpha \ordplus  1$.

\emph{\olref{exrankpair}.} Berdasarkan \olref[spine][rank]{ranksupstrict}.

\emph{\olref{exrankcup}.} Berdasarkan \olref[spine][rank]{ranksupstrict}.

\emph{\olref{exranktuple}.} Berdasarkan \olref{exrankpair}, diterapkan dua
kali.

\emph{\olref{exranktimes}.} Perhatikan bahwa $A \times B \subseteq
\Pow{\Pow{A \cup B}}$, lalu gunakan \olref{exrankpow} dua kali bersama
\olref{exrankcup}.

\emph{\olref{exrankunion}.} Jika $\alpha = \gamma\ordplus 1$, terdapat
suatu $c \in A$ dengan $\setrank{c} = \gamma$, dan tidak ada !!{element}
dari $A$ yang berperingkat lebih tinggi; jadi $\setrank{\bigcup A} =
\gamma$. Jika $\alpha$ merupakan ordinal limit, maka $A$ mempunyai
!!{element}s dengan peringkat yang sedekat apa pun dengan $\alpha$ (tetapi
tetap lebih kecil darinya). Karena itu, $\bigcup A$ juga mempunyai
!!{element}s dengan peringkat yang sedekat apa pun dengan $\alpha$ (tetapi
tetap lebih kecil darinya), sehingga $\setrank{\bigcup A} = \alpha$.
\end{proof}
\noindent
Kita serahkan sebagai latihan penjelasan mengapa \olref{exranktimes}
melibatkan suatu \emph{per}taksamaan.

\begin{prob}
Berikan himpunan $A$ dan $B$ sedemikian sehingga $\setrank{A \times B}=
\max(\setrank{A}, \setrank{B})$. Berikan himpunan $A$ dan $B$ sedemikian
sehingga $\setrank{A \times B}=\max(\setrank{A}, \setrank{B}) \ordplus
2$. Apakah ada peringkat lain yang mungkin?
\end{prob}

Sekarang kita juga dapat menunjukkan bahwa beberapa pengertian wajar tentang
apa artinya menyebut suatu ordinal ``hingga'' atau ``tak hingga'' ternyata
bertepatan:

\begin{lem}\ollabel{ordinfinitycharacter}
Untuk sebarang ordinal $\alpha$, pernyataan-pernyataan berikut ekuivalen:
\begin{enumerate}
	\item\ollabel{ord:notinomega} $\alpha\notin \omega$, yakni $\alpha$
	bukan bilangan asli
	\item\ollabel{ord:omegaplus} $\omega \leq \alpha$
	\item\ollabel{ord:oneplus} $1 \ordplus  \alpha = \alpha$
	%\alpha \approx \alpha \ordplus 1$
	\item\ollabel{ord:plusone} $\alpha \approx \alpha\ordplus 1$,
	yakni $\alpha$ dan $\alpha\ordplus 1$ sama banyak
	\item\ollabel{ord:infinite} $\alpha$ bersifat tak hingga Dedekind
	\end{enumerate}
\end{lem}
\noindent
Jadi, kita mempunyai lima cara yang terbukti ekuivalen untuk memahami syarat
agar suatu ordinal bersifat (tak) hingga.

\begin{proof}
\emph{\olref{ord:notinomega} $\Rightarrow$ \olref{ord:omegaplus}.}
Berdasarkan Trikotomi.

\emph{\olref{ord:omegaplus} $\Rightarrow$ \olref{ord:oneplus}.} Tetapkan
$\alpha \geq \omega$. Berdasarkan Induksi Transfinit, terdapat suatu ordinal
terkecil $\gamma$ (mungkin $0$) sedemikian sehingga terdapat ordinal limit
$\beta$ dengan $\alpha = \beta \ordplus \gamma$. Sekarang:
\[
	1 \ordplus \alpha =
	1 \ordplus (\beta \ordplus \gamma) =
	(1 \ordplus \beta) \ordplus \gamma =
	\supstrict_{\delta < \beta} (1 \ordplus  \delta) \ordplus  \gamma =
	\beta \ordplus  \gamma =
	\alpha.
\]
\emph{\olref{ord:oneplus} $\Rightarrow$ \olref{ord:plusone}.} Jelas
terdapat !!a{bijection} $f \colon (\alpha \disjointsum 1) \to (1
\disjointsum \alpha)$. Jika $1 \ordplus \alpha = \alpha$, terdapat suatu
isomorfisme $g \colon (1 \disjointsum \alpha) \to \alpha$. Sekarang
perhatikan $\comp{f}{g}$.

\emph{\olref{ord:plusone} $\Rightarrow$ \olref{ord:infinite}.} Jika
$\alpha \approx \alpha \ordplus 1$, terdapat !!a{bijection} $f \colon
(\alpha \disjointsum 1) \to \alpha$. Definisikan $g(\gamma) = f(\gamma,
0)$ bagi setiap $\gamma < \alpha$; !!{injection} ini menyaksikan bahwa
$\alpha$ bersifat tak hingga Dedekind, sebab $f(0,1) \in \alpha \setminus
\ran{g}$.

\emph{\olref{ord:infinite} $\Rightarrow$ \olref{ord:notinomega}.} Ini
merupakan \olref[z][infinity-again]{naturalnumbersarentinfinite}.
\end{proof}

\end{document}
